\section{Dos Scaling Paradigm y dos rutas de investigación de producto}

El ponente distingue primero entre next-token prediction y RL on chain of thought, y después entre «una interfaz conocida que aloja una capacidad desconocida» y «primero la convicción de producto, después entrenar al modelo para lograrla». Estos dos ejes son muy importantes: el primero responde cómo se obtiene la capacidad y el segundo, cómo entra la capacidad en el producto. Si se confunden, es fácil creer por error que un preentrenamiento más potente producirá automáticamente el comportamiento de producto correcto.

\subsection{Paradigm 1: Next-token prediction es una world-building machine}

El objetivo más común del preentrenamiento autorregresivo es
\[
\mathcal{L}_{\text{NLL}}(\theta)
=-\sum_{t=1}^{T}\log p_{\theta}(x_t\mid x_{<t}).
\]
$x_t$ es el token actual, $x_{<t}$ es el contexto previo y $\theta$ son los parámetros del modelo. Un objetivo unificado permite que el modelo absorba lenguaje, conocimiento, estilo y parte de los patrones de razonamiento; por eso el ponente lo llama world-building machine. Sin embargo, lo que recompensa es la predicción local: no define directamente cuándo termina una tarea larga ni aporta señales sobre el éxito de herramientas externas, la satisfacción del usuario o la coherencia a largo plazo.

\lecturefigure{slide-11-next-token-scaling.jpg}{Paradigm 1: next-token prediction construye un modelo del mundo en múltiples dominios}{00:07:58--00:08:39}

\begin{knowledgebox}{Por qué se amplifican los errores tempranos de token}
La distribución de generación puede escribirse como $p(x_{1:T})=\prod_t p(x_t\mid x_{<t})$. En cuanto el modelo muestrea un token que se desvía del objetivo, las distribuciones condicionales siguientes se construyen sobre una historia nueva; por ello, los personajes, los hechos y los planes de un texto largo pueden desviarse gradualmente. Este problema no significa que lo autorregresivo fracase necesariamente, sino que muestra que «alta verosimilitud local» y «éxito en tareas de largo plazo» no son la misma métrica.
\end{knowledgebox}

\subsection{Paradigm 2: RL on CoT incorpora al objetivo el resultado de las tareas largas}

El aprendizaje por refuerzo denota una trayectoria de razonamiento o de acción como $\tau=(s_0,a_0,\dots,s_T)$ y optimiza
\[
J(\theta)=\mathbb{E}_{\tau\sim\pi_{\theta}}[R(\tau)].
\]
$\pi_\theta$ es la política y $R(\tau)$ es la recompensa de toda la trayectoria. Si el reward proviene de pruebas de código, verificadores matemáticos o resultados de herramientas, el modelo obtiene una señal a nivel de tarea que va más allá de la imitación token por token. En la práctica suele añadirse una restricción respecto de una política de referencia, para evitar que la política se desvíe en exceso por perseguir la recompensa:
\[
J_{\text{reg}}(\theta)=\mathbb{E}[R(\tau)]
-\beta\,D_{\mathrm{KL}}(\pi_{\theta}\Vert\pi_{\mathrm{ref}}).
\]
$\beta$ controla el equilibrio entre mejorar el comportamiento y conservar la distribución.

\lecturefigure{slide-12-rl-cot-scaling.jpg}{Paradigm 2: optimizar tareas complejas con RL sobre trayectorias de Chain-of-Thought}{00:08:39--00:09:03}

\begin{warningbox}{CoT no es automáticamente un «proceso interno» verídico}
Chain-of-Thought puede ofrecer trayectorias de cómputo más largas, admitir llamadas a herramientas y verificaciones intermedias, pero que el texto sea legible no garantiza que sea faithful. El modelo puede dar una respuesta correcta acompañada de una explicación errónea, o aprender como fórmula el formato que prefiere la recompensa. La validación debe apoyarse primero en el outcome externo, en verificaciones ejecutables y en pruebas contrafactuales, en lugar de recompensar solo lo que «parece razonamiento».
\end{warningbox}

\lecturefigure{slide-13-create-build-make.jpg}{Create, build, and make: la ampliación de capacidades debe concretarse en obras que puedan terminarse}{00:09:03--00:09:35}

\teachervoice{00:07:58--00:09:34, en clase se llama a next-token prediction y a RL on CoT dos scaling paradigm. Lo que el docente enfatiza es el cambio en la señal de supervisión: el primero aprende una distribución y el segundo puede incorporar al entrenamiento la retroalimentación sobre el resultado de tareas complejas; esto no significa que todas las tareas largas ya estén resueltas.}

\subsection{Dos rutas: Capability-first y Product-belief-first}

La primera ruta parte de una capacidad nueva y busca un form factor que el usuario ya conozca; la segunda plantea primero una convicción de producto y después construye datos, eval y entornos para que el modelo muestre ese comportamiento. La primera sirve para explorar «qué aprendió a hacer de pronto el modelo»; la segunda, para buscar una experiencia de producto coherente a largo plazo. Los proyectos reales suelen ir y venir entre ambas.

\lecturefigure{slide-14-two-product-research-paths.jpg}{Las dos rutas del research-driven product}{00:09:34--00:10:00}

\lecturefigure{slide-15-familiar-form-unfamiliar-capability.jpg}{Ruta uno: alojar una capacidad desconocida en un form factor conocido}{00:09:43--00:10:06}

\begin{importantbox}{El triángulo Capability, Affordance y Eval}
\term{Capability} es lo que el modelo puede hacer potencialmente; \term{affordance} es cómo la interfaz permite al usuario invocar, editar y deshacer; \term{eval} es la evidencia con la que el sistema juzga que algo está terminado. Los fracasos de producto suelen deberse a un desajuste entre los tres: el modelo sabe hacerlo pero la interfaz no puede expresarlo, la interfaz puede activarlo pero el eval no mide el resultado real, o el eval optimiza una métrica sustituta distinta del objetivo del usuario.
\end{importantbox}

\subsection{Caso uno: CLIP fashion search}

CLIP proyecta imágenes y textos en un espacio de embedding compartido. Si la codificación de la imagen es $f_I(i)$ y la del texto es $f_T(q)$, la puntuación de recuperación suele usar la similitud coseno
\[
s(i,q)=\frac{f_I(i)^{\top}f_T(q)}{\lVert f_I(i)\rVert\lVert f_T(q)\rVert}.
\]
Esta capacidad desconocida se colocó dentro de una interfaz de búsqueda conocida: el usuario da un texto o una prenda objetivo y el sistema devuelve resultados visualmente similares. Al leer los tres estados consecutivos, conviene comparar cómo la intención de la query cambia el conjunto de resultados, en lugar de tomar cualquier grupo de imágenes atractivas como una búsqueda exitosa.

\lecturefigure{slide-16-clip-fashion-target.jpg}{Fashion search: la recuperación parte de la prenda objetivo y sus atributos visuales}{00:10:06--00:10:11}

\lecturefigure{slide-17-clip-fashion-yellow-search.jpg}{Fashion search: la semántica del color amarillo y del modelo de prenda cambia la distribución de resultados}{00:10:11--00:10:23}

\lecturefigure{slide-18-clip-fashion-style-search.jpg}{Fashion search: una consulta de estilo en lenguaje natural devuelve un nuevo conjunto de candidatos}{00:10:23--00:10:54}

\begin{knowledgebox}{Lectura de la figura: una demo de recuperación debe medir al menos cuatro cosas}
Primero, relevance: si los resultados corresponden a la query; segundo, diversity: si solo devuelve casi duplicados; tercero, controllability: si modificar el color, el material o la silueta cambia los resultados de forma local; cuarto, failure visibility: si el sistema permite al usuario ver por qué hay coincidencia. El embedding compartido aporta el mecanismo de ordenamiento, pero no garantiza automáticamente equidad, disponibilidad en inventario ni que los términos de estilo signifiquen lo mismo en distintas culturas.
\end{knowledgebox}

\subsection{Caso dos: 100K context y carga de archivos}

Cuando la capacidad de contexto largo entra en el producto, el form factor más intuitivo es «subir un archivo y hacer preguntas». Sin embargo, que la token window sea lo bastante grande solo indica que la entrada puede recibirse, no que el modelo haya encontrado las páginas clave, agregado correctamente las cifras o mantenido la correspondencia con las fuentes. Una tarea de análisis de un 10-K de 85 páginas debe validar al menos cuatro tipos de error: retrieval, calculation, citation y omission.

\lecturefigure{slide-19-100k-context-business-analyst.jpg}{La promesa de producto de 100K context: que Claude actúe como analista de negocios}{00:10:54--00:11:00}

\lecturefigure{slide-20-100k-context-input.jpg}{Lado de entrada: un Form 10-K de 85 páginas y una solicitud de análisis concreta}{00:11:00--00:11:13}

\lecturefigure{slide-21-100k-context-analysis.jpg}{Lado de salida: los indicadores financieros clave, las tablas y las explicaciones se organizan como análisis}{00:11:13--00:11:34}

\begin{warningbox}{Long context no equivale a long-context reliability}
La ventana de contexto es un límite de capacidad; la confiabilidad depende además del sesgo de posición, los cálculos entre páginas, el análisis de tablas, la evidencia contradictoria y las citas. Si el producto solo prueba «si se puede subir», el modelo puede entregar, justo en el caso más peligroso, un resumen fluido que omite pasivos clave. Deben construirse cuatro grupos de pruebas: respuestas localizables, agregación entre secciones, ausencia de respuesta y archivos contradictorios.
\end{warningbox}

\subsection{Caso tres: Summarizer e interfaces de incertidumbre}

La demo de Summarizer concentra en una interfaz de documento los párrafos repetidos, la calibración y la edición. La convicción de producto que expresa no es «cuanto más corto el resumen, mejor», sino permitir que el usuario vea qué oraciones se conservaron, qué contenido se repite y qué conclusiones requieren nueva verificación. Si el sistema puede generar $K$ candidatos, la coherencia entre candidatos puede usarse como indicio de incertidumbre, pero un error coherente no debe tomarse como confianza. En comparación con el caso anterior de 100K context, aquí importa más cómo el usuario revisa y modifica el resultado: el contexto largo resuelve «cuánto puede leer», y la summarizer surface resuelve «cómo comprimir la evidencia en un artifact que siga siendo rastreable».

\lecturefigure{slide-22-o1-summarizer-calibration.jpg}{Summarizer: detección de repeticiones, P(IK) e interfaz de documento editable}{00:12:41--00:13:39}

\begin{knowledgebox}{De la interfaz de producto a las muestras de evaluación}
Un flujo de trabajo real puede registrar como evidence estructurada las aceptaciones, modificaciones, eliminaciones, reintentos y saltos a citas del usuario. El evento mínimo puede escribirse como
\[
e=(\text{task},\text{model version},\text{action},\text{artifact diff},\text{outcome}).
\]
Solo cuando el outcome está definido pueden los registros convertirse en datos de entrenamiento o de eval; recopilar únicamente «dónde hizo clic el usuario» lleva fácilmente a confundir desconcierto con preferencia.
\end{knowledgebox}

\begin{lstlisting}[caption={Ciclo mínimo de validación de un prototipo Capability-first}]
task = observe_real_user_job()
prototype = expose_capability_in_familiar_surface(task)
trace = run_with_permissions_and_logging(prototype)
failures = classify(trace, by=[outcome, control, trust, cost])
eval_set = convert_failures_to_reproducible_cases(failures)
ship_only_if(eval_set, product_constraints).passes()
\end{lstlisting}

\teachervoice{00:09:34--00:13:39, el ponente usa Inter Alia, 100K context y summarizer para ilustrar la primera ruta: primero aparece una capacidad desconocida y después se busca un form factor conocido. La lección práctica es que la interfaz debe permitir al usuario entender, controlar y verificar la capacidad, en lugar de limitarse a conectar el modelo a un cuadro de entrada.}

\subsection{Resumen de la sección}

Next-token prediction aporta un conocimiento previo amplio, y RL on CoT puede incorporar al objetivo el outcome de las trayectorias; ninguno de los dos determina por sí solo el producto. La ruta Capability-first necesita un form factor adecuado, y la ruta product-belief-first necesita traducir la convicción en entrenamiento y evaluación. La lección común de CLIP, el contexto largo y summarizer es que la capa de producto debe definir de forma explícita la controlabilidad y la evidencia.
